\section{Usuarios}

\subsection{Establecemos la contraseña de \textit{root}}
\begin{frame}
\begin{tcolorbox}[title=Establecemos la contraseña del usuario \textit{root}]
    \orden{passwd}
\end{tcolorbox}
\end{frame}


\subsection{Nuevo usuario \textit{non-root}}
\begin{frame}
    Seguiremos los siguientes pasos para la creación de nuestro usuario {\color{red}\textbf{viriato}} con su contraseña {\color{red}\textbf{iberia}}.
      \begin{tcolorbox}[title=Creamos el grupo de \textbf{viriato}.]
       	\orden{groupadd {\color{red}viriato}}
      \end{tcolorbox}\pause
   	
    \begin{tcolorbox}[title=Creamos el usuario.]
   	\orden{useradd {\color{blue}viriato} {\color{brown}-m} {\color{red}-g viriato}}
    \tcblower
    \scriptsize
   	Con el parámetro {\color{brown}\textbf{-m}} indicamos que debe crear el directorio {\ttfamily\color{red}/home/viriato} y {\color{green}\textbf{-g}} para indicar el grupo principal del usuario.
    \end{tcolorbox}\pause
   	
    \begin{tcolorbox}[title=Le ponemos al usuario una contraseña.]
     \orden{passwd {\color{blue}viriato}}
    \end{tcolorbox}
\end{frame}

 \subsection{Privilegios \textit{sudo} para \color{blue}viriato}
\begin{frame}
		Queremos darle a {\color{brown}viriato} privilegios de tipo \textit{sudo}.

    \begin{tcolorbox}[title=Para ellos debemos editar el fichero {\ttfamily\color{red}/etc/sudoers}]
		\orden{nano /etc/sudoers}
    \end{tcolorbox}
		
		
   \begin{tcolorbox}[title=Y añadimos la siguiente línea]
   	 \orden{viriato ALL=(ALL:ALL) NOPASSWD:ALL}
     \tcblower
     \scriptsize
 	   Guardamos y cerramos el fichero.
   \end{tcolorbox}
		
\end{frame}
